\documentclass{article}
\usepackage[utf8]{inputenc}
\usepackage[T1]{fontenc}
\usepackage[margin=2cm]{geometry}
\usepackage{hyperref}
\sloppy
\begin{document}
\section*{Historical source}
Encyclopaedia Britannica scan, source row 1614. Parse 3.5 extracted the text. Super 3.5 VL corrected it from an enlarged crop of each block. It can still contain recognition errors.
Page headers, footers, and page numbers are removed.

\subsection*{Untitled continuation (E0133)}
\textit{Continues an entry from an earlier page.}

shire, situated on the river Swale, fifteen miles north-
west of York. It sends two members to parliament.

\subsection*{THIRSTI (E0134)}
THIRSTI, an overall definition, parlament a deficiency in the false to moetles the inward parts of the mouth; hence arise after Hong before for think : it is a lymptom generally attending benefit afforded.

\subsection*{THISTLE (E0135)}
THISTLE, in botany. See Carduus.

Order of the Thistle, or of St Andrew, a military order of knighthood in Scotland, the rise and infiltration whereof is variously related by different authors. Leley, bishop of Roos, reports, that the night before the battle between Athelstan king of Northumberland, and Hungus king of the Picts, a bright cross, in form of that whereon St Andrew (the tutelar saint of Scotland) suffered martyrdom, appeared to Hungus; who having gained the victory, ever after bore the figure of that cross on his banners. Others assert, that Achaius king of Scotland first instituted this order, after having made the famous league offensive and defensive with Charlemagne king of France. But although the thistle had been acknowledged as the symbol of the kingdom of Scotland from the reign of Achaius, yet some refer the beginning of this order to Charles VII, of France. Others place the foundation of it as low as the year 1500.

The chief and principal ensign is a gold collar composed of thistles and sprigs of rue interlinked with amulets of gold, having pendant thereunto the image of St Andrew with his cross, and the motto, Nemo me impune lacessit.

The ordinary or common ensign worn by the knights, is a star of four silver points, and over them a green circle, bordered and lettered with gold, containing the said motto, and in the centre is a thistle proper; all which is embroidered on their left breast, and worn with the collar, with a green ribband over the left shoulder, and brought under the right arm; pendant thereto is the image of St Andrew, with his cross, in a purple robe, within an oval of gold enamelled vert, with the former motto; but sometimes they wear, encircled in the same manner, a thistle crowned.

About the time of the reformation, this order was dropped, till James II. of England refumed it, by creating eight knights : however, the revolution unfettled it again, and it lay neglected till queen Anne, in 1703, refored it to the primitive design, of twelve knights of St Andrew.

\subsection*{THLASPI (E0136)}
THLASPI, in botany, a genus of the tetradynamia filiculosa class. The pod is emarginated, heart-shaped, containing many feeds. There are ten species, six of them natives of Britain, viz. the campestre, or mithridate; the arvense, or treacle-muffard; the perfoliatum, or perfoliate treacle-muffard; the frutum, or perennial mithridate muffard; the montanum, or mountain mithridate-muffard; and the bursa pastoris, or shepherd's purse.

\subsection*{THOMANS (E0137)}
THOMANS, Thomists, or Christians of St Thomas, a people of the East-Indies, who, according to the tradition, received the Gospel from St Thomas. Upon the arrival of the Portuguese at Calicut, in their first voyage to the Indies, they met with ancient Christians, who pretended to be descended from those converted by St Thomas. The Thomans being informed of a new people arrived among them, who bore a particular veneration for the

crois, fet entheaffors to them, to make an alliance with them, and to fulfill their affiliates against the Granite prince, by whom they were greatly profitable. A mixture of opinions, with a total interruption of pallors, femonines for travel years together, occasional that horrible chain that religion was in at the interval of the Portuguese; for a specimen whereof we shall add their names of celebration the exclusion. Over their altar was a kind of gallery, and while the pristine living being beginning of the office below, a cake of floor of rice was trying to oil, or butter, above , when enough, the cake was let down in a bullet upon the slary, where target conferred it: as to the other few pieces, for wine they died a kind of body or struck, various prepared in that country. Now their ordination much more regular; the archdeacon, who was fomentes more repled than the baby handful, frequently obtained prices their other album went enntric. The Portuguese, for these two latl centuries, have laboured the reforma tion of the church, and have employed both the ecolela fific and leaching them better; for this end they have called the Thomas bibops to the council at Goa, have inherited, charged them, for even for them for induction on Portugal and Rome? He finding that they were fill apt to relale at their return, and that no good was like to be done with them, they refused to exclude them once for all, and, and to spoon in European bibings in their room. These proceedings have rendered the Portuguese familyely close to the Thomas.

St THOMAS, a city of the hither India, on the coast of Coromandel, three miles south of Fort St George; subject to the Portuguese.

St Thomas is also an island in the Atlantic ocean, situated under the equator, in 8° E. long.

St Thomas is also a town of Guiana, in South America, situated on the river Oronoko; subject to Spain.

St. Thomas's Day, a festival of the Christian church, observed on Dec. 21. in commemoration of St Thomas the apostle.

St Thomas of Canterbury's day, a festival of the Romish church, observed on Dec. 29, in memory of Thomas Becket archbishop of Canterbury, who was murdered, or, as the Romanists say, martyred, in the reign of king Henry II.

\subsection*{THOMISM (E0138)}
THOMISM, the doctrine of St Thomas Aquinas, and his followers the Thomists, chiefly with regard to predestination and grace. There is some doubt what the true genuine thomism is, but there are authors who distinguish the thomism of St Thomas from that of the dominicans. Others again make thomism no other than a kind of jansenism disguised: but jansenism, it is known, has been condemned by the popes, which pure thomism never was: in effect, the writings of Alvarez and Lemos, who were appointed by their order to lay down and defend, before the holy see, the dogmata of their school, have since been reputed the rule of pure thomism.

\subsection*{THOMISTS (E0139)}
THOMISTS, a sect of school-divines, who maintain thom-
mism. See the preceding article.

mim. See the preceding article.
THORACIC duct. See Anatomy, p. 282.
THORAX. in anatomy. See Anatomy, p.

\subsection*{THORAX (E0140)}
THORAX, in anatomy. See Anatomy, p. 277.

\subsection*{THORAX (E0141)}
THORAX, in anatomy. See ANATOMY, p. 277.
THORN, a city of Poland, in the province of royal Prussia, situated on the river Vistula: E. long. 19°, and N. lat. 52° 40'.


\section*{Source record}
\begin{description}
\item[Scan ID] \texttt{britannica:309f765b1ecc83e0a7f7762e44b4a927a7555443:train:1614}
\item[Image file] \texttt{page-1614-ea352b487d65.jpg}
\item[Image SHA-256] \texttt{ea352b487d650c713a3e6bdd8e4bc1ded85e0d7558c3fa1d9cdbce89abc38ca3}
\item[Dataset] NationalLibraryOfScotland/encyclopaedia\_britannica\_illustrated at revision \texttt{309f765b1ecc83e0a7f7762e44b4a927a7555443} (CC0-1.0)
\item[Extraction] model \texttt{nvidia/nemotron-3.5-parse}, raw reply \texttt{extract/raw/page-1614-ea352b487d65.json}
\item[Raw reply SHA-256] \texttt{10a19571244416cc047af8fea843ac2265ad882874221815ef167d9815a00f9a}
\item[Correction] model \texttt{nvidia/Nemotron-3.5-Super-VL}, block crops, raw replies \texttt{correct/raw/crop\_nothink/page-1614-ea352b487d65-b*.json}; parse: 1, parse\_fallback: 1, super\_vl: 19
\item[Running header] T H O ( 896 )
\end{description}
\end{document}
